% =====================================================================
% 2. DATA DESCRIPTION
% =====================================================================
\section{Data description}

\subsection{Data bases}

The data are the \emph{Glioma Grading Clinical and Mutation Features} dataset of the UCI Machine
Learning Repository \cite{tasci2022dataset}, built from the TCGA lower-grade glioma and
glioblastoma projects \cite{tcga2015lgg,brennan2013} and introduced in \cite{tasci2022}. The original
export contains \NRaw{} patients; the encoded UCI version contains \NClean{}. Each patient is
described by age at diagnosis, sex, race, the somatic mutation status (mutated or not) of 20 genes,
the TCGA project (TCGA-LGG or TCGA-GBM), the histological diagnosis and the grade. The final dataset
has \NLGG{} LGG and \NGBM{} GBM patients (\PctGBM\% GBM).

Three properties of these data limit what any model can learn, and they shape the critical analysis
of Section~\ref{sec:prdl-results}:
\begin{itemize}
  \item \textbf{Labels from older criteria.} The cohorts were collected over many years
        (TCGA-GBM between 1989 and 2011, with 76\% of patients diagnosed in or after 2002
        \cite{zakharova2023}) and graded before the 2021 WHO classification. TCGA glioma profiles
        have since been re-annotated following the 2021 criteria \cite{zakharova2023,leiriademendonca2025}%
        \todoverify{quote from 41.pdf or 42.pdf stating how many TCGA cases change category when the
        original labels are replaced by WHO 2021 labels}.
        The grade used here as ground truth is therefore partly outdated: under the current criteria,
        some tumours labelled LGG would be classified differently.
  \item \textbf{Point mutations only.} Several markers that now define the type or the grade are
        copy-number or promoter alterations absent from the dataset: the 1p/19q codeletion
        (oligodendroglioma) \cite{louis2021}; the homozygous deletion of CDKN2A/B, ``a molecular marker
        of WHO grade 4 in IDH-mutant astrocytomas''; and the TERT promoter mutation, EGFR amplification
        or the +7/$-$10 signature, which define glioblastoma \cite{weller2021}. These alterations are
        detected with copy-number assays such as FISH \cite{sahm2023}. A model can only approximate them
        through correlated point mutations: in the TCGA lower-grade gliomas, the 1p/19q-codeleted tumours
        ``harbored mutations in CIC, FUBP1, NOTCH1, and the TERT promoter'' \cite{tcga2015lgg}.
  \item \textbf{A homogeneous cohort.} \PctWhite\% of the patients are white (\NWhite{} of
        \NClean); only \NBlack{} are Black or African American, \NAsian{} Asian and \NNative{}
        American Indian or Alaska Native. Any conclusion about race, and the transfer of the models
        to other populations, is unsupported.
\end{itemize}

\subsection{Data preparation}

The preparation followed three principles, each of which can be justified from theory.

\textbf{Reproducibility of the published cleaning.} Rather than loading the UCI file, the clean
dataset was rebuilt from the original export, keeping the patient identifier and the diagnosis,
and checked value by value against the UCI file. The \NRemoved{} excluded patients
(\NRemovedAllMissing{} with every clinical field missing, \NRemovedRace{} with race not reported and
\NRemovedAge{} with missing age) are split almost evenly between classes (\NRemovedLGG{} LGG,
\NRemovedGBM{} GBM), so complete-case analysis does not distort the class balance.

\textbf{No target leakage.} Leakage occurs when information that would not be available at
prediction time, or that encodes the target, enters the model \cite{kaufman2012}. Two columns do
exactly that: the TCGA project coincides with the grade for all \NProjectGrade{} patients, and the
histological diagnosis determines it (every glioblastoma is GBM and every other diagnosis is LGG).
Both are excluded from all models. The EDA (Section~\ref{sec:prdl-eda}) quantifies the effect of
ignoring this rule.

\textbf{Correct variable types.} Race is a nominal variable whose integer codes have no order, so
it is one-hot encoded; treating the codes as a number would impose an arbitrary linear effect.
Mutations are binary and are kept as 0/1 indicators. All transformations that learn from data
(standardisation, encoding) are fitted inside the cross-validation (Section~\ref{sec:prdl-setup}).
